\section{Extended motivation and design history}
\label{sec:motivation}

This appendix expands \cref{sec:main-introduction} with the external evidence,
the design decisions that were taken and retired, and the reasoning behind the
registered positive condition.  It is not a second statement of the paper's
contributions: those are stated once, in the body, and this section is
subordinate to them wherever the two could be read differently.

Runtime evaluation of agents already shows the gap this paper is built to
measure.  In one batch of tasks, a single lifecycle-wide harness patch rescued
six baseline failures and regressed one baseline success in the same batch;
across a frontier model's 39 valid patches, 21 reduced batch success; and the
worst affected batch fell from 7/10 to 0/10 \citep{harnessr12026}.  Their own
diagnosis is blunt: ``a plausible diagnosis can still compile into overly
aggressive runtime behavior.''  This motivates but does not identify the
checkpoint-level regimes we define: a single flipped rollout is one realization
of a stochastic process, not an estimate of the effect of intervening at that
checkpoint, and we do not treat it as one.

Those observations are someone else's.  Our initial design simulation found
that targeting a semantic fault at checkpoints the registered risk score places
low raises apparent power sharply.  That result is now treated as a warning,
not an operating point: score-conditioned placement constructs an anti-ordered
population using the score under test.  Prediction-conditioned selection also
changes which outcomes later become visible in deployed online systems, an
adjacent feedback problem documented as a winner's-curse effect
\citep{staleness2026}.  The primary positive condition is
therefore a randomized, risk-blind feature-engineering--configuration
interaction, while the targeted injector is retained only as a historical
adversarial-stress calculation (\cref{sec:experiments,app:faults}).  These are
design decisions; no result from ongoing execution is admitted here.

The question is simple.  An autonomous experiment agent periodically reaches a
checkpoint at which someone must decide whether to let it continue, verify its
work, or recover it.  Predicting whether the checkpoint will fail is one
question; measuring whether intervening would have helped is a different one.
We ask when the first suffices to answer the second, and how to act on the
answer without making a running system worse.  Both questions are stated here
as precise, testable claims with an instrument that would settle them.  With no
qualifying result admitted, they remain predictions rather than findings.

The observation that prediction and intervention value differ is established
\citep{interventionparadox2026,calibrationnotcontrol2026}; related production
systems also report that better offline prediction or supervision need not
translate proportionally into better constrained placement, serving, or
orchestration decisions \citep{fleda2026,staleness2026,ajs2026,gstar2026}.
We take the separation as the starting point rather than the contribution.  The
primary research question is
prospective and causal: within a registered population of eligible checkpoints,
\emph{when} do the two come apart, and can a one-step policy exploit the
difference to select among feasible regimes better than the strongest
constrained non-personalized policy?  ``Better'' means higher independently
validated completion probability at a registered prevalence, subject to
no-fault noninferiority, invalid-output tolerances, and resource constraints.
These are planned questions, not findings.

The theoretical results are individually elementary and are presented as such in
\cref{sec:theory}.  What the one-step restriction buys is that their hypotheses
become properties of the design rather than assumptions; the contribution the
body claims is the measurement contract of \cref{sec:main-scope}, not the
arguments themselves.
